%% FORMLÆRE: LÆREBOK — Kapittel 6–9 + Appendiks
%% Denne fila vert \input{} i hovuddokumentet.
%% Iver Raknes Finne, Institutt for design, AHO
%% Utkast 01 — 2026-03-30
%%
%% Avhengig av teorem-miljø, kommandoar og pakkar definerte i hovudfila.
%% Krev: amsmath, amssymb, amsthm, bm, booktabs, longtable, enumitem,
%%        hyperref, cleveref, og dei eigendefinerte \MC, \Nav, etc.


%% =====================================================================
%%  KAPITTEL 6: SUBSTRAT-UAVHENGIGHEIT OG EMERGENS
%% =====================================================================
\chapter{Substrat-uavhengigheit og emergens}
\label{chap:substrat}

\begin{quote}
\itshape
Kvar gong du ser ein termostat justere temperaturen, ein algoritme
optimere ei form, eller ein organisme regenerere eit lem, ser du
den same strukturen: eit system som har eit mål, kan registrere
avstand frå målet, og kan justere seg mot det. Materialet varierer.
Logikken er identisk.
\end{quote}

\section{Navigasjon er ein eigenskap ved organisering}

I kapittel~5 definerte vi navigatoren som eit trippel
$\Nav = (G, \mu, \alpha)$: eit målsett, ein avstandsfunksjon,
og eit styringsfelt. Vi stilte tre vilkår --- (a)~systemet har ein
måltilstand, (b)~det kan registrere avstand frå målet, og (c)~det
kan justere seg mot målet i gjennomsnitt --- men vi sa ingenting
om \emph{kva} systemet er laga av.

Ein termostat oppfyller alle tre vilkåra. Det same gjer ein
gradient-descent-algoritme, ein koloni av celler som byggjer eit
organ, og ein handverkar som dreier ein stolbein. Substratet --- metall,
silisium, organisk vev, menneskeleg medvit --- varierer radikalt.
Den formelle strukturen gjer det ikkje.

Denne observasjonen er ikkje triviell. Ho seier at navigasjon er
ein eigenskap ved \emph{korleis} eit system er organisert, ikkje
ved \emph{kva} det er laga av. Konsekvensen er djuptgripande:
dersom vi vil forstå kvifor former ser ut som dei gjer, treng vi
ikkje avgrense oss til éin type formgjevar. Vi kan samanlikne
celler og algoritmar, handverkarar og marknader, med same
vokabular.


\section{Postulat: Substrat-uavhengig navigasjon}

\begin{postmark}[Substrat-uavhengig navigasjon]\label{post:L.7.1}
Vilkåra (a), (b) og (c) i definisjonen av ein navigator kan
oppfyllast av eitkvart fysisk system --- biologisk, algoritmisk,
mekanisk eller distribuert --- som:
\begin{enumerate}[label=(\roman*),nosep]
  \item har tilstandar (posisjonar i eit konfigurasjonsrom),
  \item kan registrere avstand frå ein målkonfigurasjon, og
  \item kan justere tilstanden sin som respons.
\end{enumerate}
Navigasjon er ein eigenskap ved den \emph{funksjonelle
organiseringa} til systemet, ikkje ved materialet det er laga av.
\end{postmark}

\paragraph{Falsifiseringsvilkår.}
Postulatet fell dersom det kan visast at minst eitt av vilkåra
(a)--(c) \emph{nødvendigvis} krev eit spesifikt fysisk substrat
--- til dømes at organisk materie er naudsynt. Meir presist:
dersom det finst formgjevingsprosessar der ein menneskeleg
navigator konsekvent produserer former som ingen algoritmisk
navigator kan produsere, \emph{ikkje} på grunn av reknekapasitet
eller sensorisk oppløysing, men fordi sjølve navigasjonen krev
organisk substrat, då fell postulatet.

\medskip
Legg merke til kor spesifikt kravet er. Det held ikkje å vise at
ein menneskeleg designar er \emph{betre} enn ein algoritme. Det
må visast at det finst ein \emph{prinsipiell} barriere, ikkje ei
praktisk avgrensing. Skilnaden mellom ``endå ikkje realisert''
og ``prinsipielt umogleg'' er skilnaden mellom ingeniørarbeid
og falsifisering.


\section{Posisjonsinvarians}

Før vi kan samanlikne navigatorar på tvers av substrat, treng vi
ei presis påstand om det rommet dei navigerer \emph{i}. Formrommet
$\MC$ vart definert i kapittel~1 som mengda av alle
eigenskapsvektorar som oppfyller dei fysiske avgrensingane for ein
funksjonell klasse. Spørsmålet er: avheng posisjonen i formrommet
av korleis objektet vart laga?

\begin{satmark}[Posisjonsinvarians]\label{thm:L.7.2}
Posisjonen $\bm{x} \in \MC$ til eit realisert objekt er bestemt
av dei målbare eigenskapane til objektet, ikkje av prosessen som
produserte det. To objekt med identiske eigenskapsvektorar
okkuperer same punkt i $\MC$, uansett om dei vart laga for hand,
av maskin, av algoritme, eller av biologisk vekst.
\end{satmark}

\begin{proof}[Bevis, steg for steg]
\begin{enumerate}
  \item \textbf{Utgangspunkt.}\;
    Eigenskapsvektoren $\bm{x}(a) = (x_1(a), \ldots, x_n(a))$ er
    definert utelukkande av dei målbare geometriske eigenskapane
    til objektet~$a$ (høgd, breidd, masse, kurvatur, symmetri, osb.).
  \item \textbf{Nøkkelobservasjon.}\;
    Desse eigenskapane er \emph{intrinsiske} for objektet. Dei er
    ikkje funksjonar av produksjonshistoria. Eit stykke tre har
    same kurvatur anten det vart forma med kniv eller med
    CNC-fres.
  \item \textbf{Konklusjon.}\;
    La $a$ vere eit handlaga objekt og $b$ eit maskinlaga objekt.
    Dersom $x_i(a) = x_i(b)$ for alle $i = 1, \ldots, n$, so er
    $\bm{x}(a) = \bm{x}(b)$, og dei to objekta okkuperer same
    punkt i~$\MC$. \qedhere
\end{enumerate}
\end{proof}

\paragraph{Kva teoremet ikkje seier.}
Posisjonsinvarians betyr \emph{ikkje} at produksjonsprosessen er
irrelevant. Prosessen avgjer kva \emph{region} av $\MC$ som er
tilgjengeleg --- substratet avgrensar trajektorien. Biologisk vekst
kan ikkje produsere rette vinklar. Dreiing kan ikkje produsere
asymmetri. Kvar slik avgrensing er ein usynleg vegg i formrommet.
Men for dei punkta som \emph{er} tilgjengelege, er posisjonen
sjølv substrat-invariant.

\paragraph{Konsekvens.}
Sidan formrommet er substrat-invariant, kan vi plassere objekt frå
\emph{alle} substrat i \emph{same} rom og samanlikne dei direkte.
Stolen forma av ein snikkar og stolen generert av ein algoritme
okkuperer begge punkt i same $\MC$. Spørsmålet er ikkje om dei
\emph{kan} samanliknast, men kva samanlikninga \emph{avslører}.


\section{Emergent navigasjon}

\begin{konmark}[Emergent navigasjon]\label{conj:L.7.3}
La $\Nav^1, \Nav^2, \ldots, \Nav^S$ vere ei samling
sub-navigatorar, kvar med sitt eige trippel
$(G_j, \mu_j, \alpha_j)$. Det kollektive systemet
$\Nav = \{\Nav^1, \ldots, \Nav^S\}$ kan realisere eit
navigator-trippel $(G, \mu, \alpha)$ der $G$, $\mu$ og $\alpha$
\emph{ikkje} er funksjonar av nokon einskild sub-navigators
trippel, men oppstår frå interaksjonsdynamikken.
\end{konmark}

\paragraph{Den formelle strukturen.}
Definer den kollektive tilstanden som
$\bm{X} = (\bm{x}^1, \ldots, \bm{x}^S) \in (\MC)^S$.
Kvar sub-navigator $\Nav^j$ genererer ein straum
$\dot{\bm{x}}^j = \alpha_j(\bm{x}^j)$ i sitt eige formrom.
Når sub-navigatorane interagerer --- gjennom delt seleksjonstrykk,
kommunikasjon eller fysisk kopling --- vert den kollektive
dynamikken:
\[
  \dot{\bm{x}}^j = \alpha_j(\bm{x}^j)
  + \sum_{l \neq j} \phi_{jl}(\bm{x}^j, \bm{x}^l)
\]
der $\phi_{jl}$ kodar interaksjonen mellom sub-navigatorane $j$
og~$l$. Det kollektive systemet har ein makroskopisk navigator
$(G, \mu, \alpha)$ dersom det finst ei grovkorning
$\Pi : (\MC)^S \to \MC$ slik at den projiserte dynamikken
$\dot{\bm{y}} = \Pi_* \dot{\bm{X}}$ oppfyller
navigatordefinisjonens tre vilkår --- og trippelet ikkje kan
dekomponerast som ein funksjon av nokon einskild
$(G_j, \mu_j, \alpha_j)$.

\paragraph{Kvifor er dette ein konjektur og ikkje eit teorem?}
Vi påstår at slik emergens \emph{finst}, men vi har ikkje
presisert vilkåra på $\phi_{jl}$ som garanterer det. Kva
koplingsstyrke trengs? Kva symmetri? Kva grenseverdiar? Desse
spørsmåla er opne. Konjekturen er ei presis formulering av eit
uløyst problem, ikkje eit ferdig bevis.

\paragraph{Empirisk grunngjeving.}
Syntetiske organismar --- xenobotar --- sette saman av
froskeembryoceller utan evolusjonshistorie, navigerer likevel mot
koherente morfologiske konfigurasjonar. Navigasjonskapasiteten er
ikkje koda i nokon einskild celles genom; ho oppstår frå
intercellulær bioelektrisk signaldynamikk. I designkonteksten
svarar dette til prosjekt der det ferdige resultatet ikkje var
spesifisert av nokon einskild deltakar, men emergerte frå
forhandlinga mellom dei.

At navigasjon kan vere \emph{emergent}, ikkje programmert, er det
sterkaste argumentet for at rammeverket gjeld utover dei enklaste
navigatorane.


\section{Designhistorisk døme: Topologi møter biologi}

I 1917 skildra D'Arcy Thompson korleis biologiske strukturar ---
beinkonstruksjonar, trebjelkar, korallreiv --- konvergerer mot
former som minimerer materialebruk under mekanisk last. Hundre år
seinare produserer topologisk optimering --- ein reint algoritmisk
prosess som iterativt fjernar materiale frå eit designdomene ---
nesten identiske former under same fysiske gradientar.

Konvergensen er slåande. Eit lårbeins trabekulære struktur og
utdata frå ein topologisk optimeringskøyring ser ut som tvillingar.
Substratet er heilt ulikt: organisk vev på den eine sida,
silisium og aluminium på den andre. Seleksjonstrykka er
dei same: mekanisk last, materialkostnad, strukturell integritet.
Resultata konvergerer fordi tilpassingslandskapet er
substrat-uavhengig.

\medskip
\begin{center}
\begin{tabular}{@{}lll@{}}
\toprule
& \textsc{Biologisk bein} & \textsc{Topologisk opt.} \\
\midrule
Substrat & Organisk vev & Metall/polymer \\
Navigator & Cellulær signalering & Algoritme (SIMP/ESO) \\
Seleksjonstrykk & Mekanisk last + metabolsk kostnad
  & Stivleik + massebudsjett \\
Resulterande form & Trabekulært nettverk & Lattice-struktur \\
Konvergens & \multicolumn{2}{c}{Same topologi under same gradientar} \\
\bottomrule
\end{tabular}
\end{center}

\medskip
Dette er posisjonsinvarians i praksis: to radikalt ulike
navigatorar, i to radikalt ulike substrat, konvergerer mot same
region av formrommet fordi dei responderer på same
seleksjonstrykk. Posisjonen i $\MC$ er bestemt av eigenskapane,
ikkje av produsenten.


\section{Gjennomarbeidd døme: Tre navigatorar, éin stol}

Tenk deg tre stolformer produserte under same funksjonelle krav
(bere 80~kg menneskeleg vekt, sitjehøgd 43--47~cm, stabil på
plant golv) og same estetiske preferanse (visuell lettheit). Tre
ulike navigatorar tek oppgåva:

\paragraph{Navigator A: Handverkaren.}
Ein snikkar med 20 års erfaring i bøk arbeider med tradisjonelle
handverktøy. Ho navigerer lokalt og inkrementelt: kvar justering
er liten, informert av taktil tilbakemelding frå materialet.
Resultatet er ein stol med lett kurva bakstykke, runde
stolbeinsnitt, og ein kompaktheit rundt 0{,}44.

\paragraph{Navigator B: Algoritmen.}
Ein topologisk optimeringsprosess med gjevne lasttilfelle og
massebudsjett, implementert i eit parametrisk CAD-system.
Navigatoren utforskar rommet sprangvis gjennom mutasjon og
seleksjon. Resultatet er ein stol med organisk, forgreina
struktur, open profil, og ein kompaktheit rundt 0{,}41.

\paragraph{Navigator C: Evolusjonen.}
Tusen år med stoltradisjon i ein kultur med same treslag og same
kroppsproporsjoner, der kvar generasjon modifiserer forgjengaren
sin stol marginalt. Navigatoren er distribuert: ingen einskild
aktør kontrollerer prosessen. Resultatet er ein vernakular
storltype med rette linjer, kraftigare dimensjonar, og ein
kompaktheit rundt 0{,}46.

\paragraph{Samanlikning.}
\begin{center}
\begin{tabular}{@{}lccc@{}}
\toprule
& \textsc{Handverkar} & \textsc{Algoritme} & \textsc{Evolusjon} \\
\midrule
Kompaktheit & 0{,}44 & 0{,}41 & 0{,}46 \\
Symmetri & 0{,}12 & 0{,}09 & 0{,}14 \\
Kurvatur & 0{,}10 & 0{,}13 & 0{,}06 \\
\midrule
Substrat & Bøk & Polymer/metall & Bøk \\
Trajektorie & Lokal, inkr. & Sprangvis & Blind, kumulativ \\
\bottomrule
\end{tabular}
\end{center}

\medskip
Alle tre stolane ligg innanfor same \emph{region} av formrommet,
fordi seleksjonstrykka er dei same. Men dei okkuperer \emph{ulike
punkt} innanfor regionen, fordi substrata (og dermed
trajektoriane) er ulike. Posisjonsinvarians garanterer at vi kan
plassere dei i same rom. Substrat-avhengig trajektorie forklarer
kvifor dei ikkje er identiske.

Legg særleg merke til algoritmen sin låge kompaktheit: utan ein
menneskeleg designars preferanse for visuelle proporsjonar, og
utan treets motstand mot visse geometriar, konvergerer algoritmen
mot ei annleis løysing innanfor same tillatne region. Substratet
formar stigen, ikkje målet.


\section{Øvingar}

\begin{exercise}[Nivå 1]\label{ex:6.1}
Ein termostat, ein gradient-descent-algoritme, og ein bie som
byggjer ein vokskake: identifiser $(G, \mu, \alpha)$ for kvar av
dei tre navigatorane. Kva er likt? Kva er ulikt?
\end{exercise}

\begin{exercise}[Nivå 1]\label{ex:6.2}
Forklar med eigne ord kvifor posisjonsinvarians \emph{ikkje} betyr
at produksjonsprosessen er irrelevant. Gje eit konkret døme der
to ulike prosessar produserer objekt på same punkt i $\MC$, og eit
døme der dei ikkje kan gjere det.
\end{exercise}

\begin{exercise}[Nivå 2]\label{ex:6.3}
Eit designteam på fem personar utviklar ein ny kontorstol. Ingen
einskild person spesifiserer den endelege forma. Bruk konjekturen
om emergent navigasjon til å modellere teamet som eit kollektiv av
sub-navigatorar. Kva er $\phi_{jl}$ i dette tilfellet? Kva
betyr det at det kollektive målet $G$ ikkje kan reduserast til
nokon av dei individuelle $G_j$?
\end{exercise}

\begin{exercise}[Nivå 2]\label{ex:6.4}
Topologisk optimering av ein bjelke under bøyelast og biologisk
beinvekst under same lasttilfelle produserer liknande strukturar.
Vel eit konkret lasttilfelle (t.d.\ cantilever-bjelke, 100~N
punktlast) og skisser korleis du ville samanlikna resultata i eit
felles formrom. Kva eigenskapar ($x_1, \ldots, x_n$) ville du
bruke?
\end{exercise}

\begin{exercise}[Nivå 3]\label{ex:6.5}
Postulatet om substrat-uavhengig navigasjon kan falsifiserast.
Design eit eksperiment som ville avgjere om ein spesifikk
formgjevingsoppgåve \emph{prinsipielt} krev organisk substrat
(ikkje berre er vanskelegare utan). Kva kontrollvariablar treng du?
Korleis ville du skilje mellom ``endå ikkje realisert'' og
``prinsipielt umogleg''?
\end{exercise}

\begin{exercise}[Nivå 3]\label{ex:6.6}
Formlær deg konjekturen om emergent navigasjon. Finn eitt
publisert døme (frå biologi, robotikk, eller designforsking) der
eit kollektivt system viser navigasjonskompetanse som ikkje kan
reduserast til nokon einskild komponent. Spesifiser $(G, \mu,
\alpha)$ for det kollektive systemet og argumenter for kvifor dei
er emergente.
\end{exercise}


%% =====================================================================
%%  KAPITTEL 7: VARIASJONSPRINSIPP
%% =====================================================================
\chapter{Variasjonsprinsipp}
\label{chap:variasjon}

\begin{quote}
\itshape
Det er ikkje posisjonen som tel. Det er stigen.
\end{quote}

\section{Stigen gjennom formrommet}

Hittil har vi skildra tilpassingslandskapet som ein overflate og
formgjeving som rørsle på denne overflata. Vi har spurt: kva
posisjonar er optimale? Men ein stol er ikkje berre ein posisjon.
Han er resultatet av ein prosess --- ein stig gjennom formrommet
frå eit utgangspunkt til eit sluttpunkt. Og prosessen
\emph{endrar} kva som er mogleg.

I dette kapittelet utvidar vi perspektivet frå posisjon til
\emph{trajektorie}. Spørsmålet er ikkje lenger ``kvar er den
beste stolen?'' men ``kva er den beste \emph{vegen} til
stolen?'' Denne reformuleringa har store konsekvensar.


\section{Stiavhengig tilpassing}

\begin{postmark}[Stiavhengig tilpassingsverdi]\label{post:L.8.1}
Tilpassingsverdien til ei realisert form avheng ikkje berre av
posisjonen hennar $\bm{x} \in \MC$, men av heile stigen
$\gamma : [0, T] \to \MC$ som førte dit. Definer det
\emph{stiavhengige tilpassingsfunksjonalet}:
\[
  S[\gamma] \;=\; \int_0^T
    g\!\left(\gamma(t),\; \dot{\gamma}(t),\; t\right) dt
\]
der $g : \MC \times \Rn \times \R \to \R$ er ein lagrangian som
avheng av posisjonen $\gamma(t)$, hastigheita
$\dot{\gamma}(t)$ (endringsraten), og tida~$t$.
\end{postmark}

\paragraph{Kva betyr dette?}
Tenk deg to identiske stolar --- same mål, same materiale, same
overflatebehandling --- som okkuperer nøyaktig same punkt i
formrommet. Den eine er originalen, teikna og laga av ein
snikkar i 1920. Den andre er ein kopi, produsert av ein fabrikk
i 2020. Dei har same posisjon i $\MC$, men ulik kulturell verdi,
ulik marknadsposisjon, og ulikt potensial til å inspirere neste
generasjon. Posisjonen åleine er utilstrekkeleg. Stigen tel.

\paragraph{Falsifiseringsvilkår.}
Postulatet fell om identiske former med ulik produksjonshistorie
konsekvent har same overlevings- og reproduksjonsrate i
formgjevingsøkologien.

\medskip
Kublers observasjon --- ``alt som vert laga er anten ein kopi eller
ein variant av noko som vart laga litt tidlegare'' --- er ei
formulering av same innsikt: posisjonar åleine er utilstrekkelege.
Vi treng \emph{former-med-historie}, ikkje berre koordinatar.


\section{Euler--Lagrange-likninga for designtrajektoriar}

Dersom realiserte designtrajektoriar er stasjonære punkt av
$S[\gamma]$ --- det vil seie at dei oppfyller $\delta S[\gamma] = 0$
--- kva for ein liknad styrer dei?

\begin{satmark}[Stasjonær stig]\label{thm:L.8.2}
Dersom $\delta S[\gamma] = 0$, oppfyller $\gamma$ Euler--Lagrange-likninga:
\[
  \frac{\partial g}{\partial \gamma}
  - \frac{d}{dt}\frac{\partial g}{\partial \dot{\gamma}} = 0.
\]
\end{satmark}

\begin{proof}[Bevis, steg for steg]
Vi fylgjer den klassiske variasjonsrekninga.

\paragraph{Steg 1: Definer ein variasjon.}
La $\gamma_\varepsilon(t) = \gamma(t) + \varepsilon \, \eta(t)$
vere ein nabostig, der $\eta(t)$ er ein vilkårleg glatt funksjon
med $\eta(0) = \eta(T) = 0$ (endepunkta er faste). Parameteren
$\varepsilon$ er eit lite reellt tal.

\paragraph{Steg 2: Formuler stasjonaritetsvilkåret.}
Vi krev at $S[\gamma]$ er stasjonær, det vil seie at retninga
til endringa er null:
\[
  \frac{d}{d\varepsilon}\bigg|_{\varepsilon=0}
  S[\gamma_\varepsilon] \;=\; 0.
\]

\paragraph{Steg 3: Rekn ut derivaten.}
Sett inn $\gamma_\varepsilon$ i $S$ og deriver:
\begin{align*}
  \frac{d}{d\varepsilon}\bigg|_{0} S[\gamma_\varepsilon]
  &= \int_0^T \!\left(
    \frac{\partial g}{\partial \gamma} \cdot \eta
    + \frac{\partial g}{\partial \dot{\gamma}} \cdot \dot{\eta}
  \right) dt.
\end{align*}
Her er $\frac{\partial g}{\partial \gamma}$ den partielle
deriverte av $g$ med omsyn til posisjonen, og
$\frac{\partial g}{\partial \dot{\gamma}}$ den partielle
deriverte med omsyn til hastigheita.

\paragraph{Steg 4: Partiell integrasjon.}
Det andre leddet inneheld $\dot{\eta}$, den deriverte av
variasjonen. Vi integrerer delvis for å flytte derivaten over på
$\frac{\partial g}{\partial \dot{\gamma}}$:
\begin{align*}
  \int_0^T \frac{\partial g}{\partial \dot{\gamma}}
    \cdot \dot{\eta}\; dt
  &= \left[\frac{\partial g}{\partial \dot{\gamma}}
    \cdot \eta \right]_0^T
    - \int_0^T \frac{d}{dt}
      \frac{\partial g}{\partial \dot{\gamma}}
      \cdot \eta \; dt.
\end{align*}

\paragraph{Steg 5: Grenseleddet forsvinn.}
Sidan $\eta(0) = \eta(T) = 0$, er klammeuttrykket null:
$\left[\frac{\partial g}{\partial \dot{\gamma}} \cdot \eta
\right]_0^T = 0$.

\paragraph{Steg 6: Sett saman.}
No har vi:
\[
  \int_0^T \!\left(
    \frac{\partial g}{\partial \gamma}
    - \frac{d}{dt}\frac{\partial g}{\partial \dot{\gamma}}
  \right) \cdot \eta \; dt \;=\; 0.
\]

\paragraph{Steg 7: Fundamentallemmaet.}
Sidan $\eta$ er \emph{vilkårleg}, er den einaste måten integralet
kan vere null for alle val av $\eta$, at det som står inne i
parentesen er null overalt:
\[
  \frac{\partial g}{\partial \gamma}
  - \frac{d}{dt}\frac{\partial g}{\partial \dot{\gamma}}
  = 0. \qedhere
\]
\end{proof}

\paragraph{Tolking.}
Euler--Lagrange-likninga seier at optimale designtrajektoriar
balanserer to krefter: $\frac{\partial g}{\partial \gamma}$, som
dreg forma mot regionar med høg tilpassingsverdi, og
$\frac{d}{dt}\frac{\partial g}{\partial \dot{\gamma}}$, som
straffar brå endringar i retning. Ein tradisjon som endrar seg
for raskt, taper akkumulert kunnskap. Ein tradisjon som aldri
endrar seg, vert fanga i eit lokalt optimum. Likninga formaliserer
balansen.


\section{Stigen modifiserer landskapet}

\begin{satmark}[Stigen endrar kva som er mogleg]\label{thm:L.8.3}
Stiavhengig tilpassing og landskapsminne (frå kapittel~4) saman
inneber at å gje form ikkje er å finne den beste \emph{posisjonen},
men den beste \emph{stigen}: prosessen modifiserer kva resultat
som er moglege.
\end{satmark}

\begin{proof}[Grunngjeving]
Frå landskapsminne (kap.~4) veit vi at kvar realisert form
$\bm{y} = \gamma(t)$ endrar dei framtidige seleksjonstrykka via
ein påverknadsfunksjon~$\eta_i$:
\[
  s_i(\bm{x},\, t + dt) \;=\;
  s_i(\bm{x},\, t) + \eta_i(\bm{y},\, \bm{x}).
\]
Tilpassingslandskapet ved tidspunkt $T$ er difor eit funksjonale
av \emph{heile} trajektorien:
\[
  f(\bm{x}, T) = \Phi\!\left(
    s_1\!\left(\bm{x},\, 0\right) +
    \int_0^T \!\eta_1(\gamma(\tau), \bm{x})\, d\tau,\;
    \ldots
  \right).
\]
To ulike trajektoriar $\gamma$ og $\gamma'$ som endar i same
punkt $\bm{x}^T$ kan stå overfor ulike framtidige landskap, ulike
tilgjengelege optima, og ulike overlevingsutsikter. Prosessen er
ikkje berre middelet til resultatet; prosessen \emph{modifiserer}
kva resultat som er moglege.
\end{proof}

\paragraph{Konsekvens for praksis.}
Ein stol laga gjennom lang handverkstradisjon har endra
designlandskapet rundt seg: han har etablert forventningar,
skapt referansepunkt, opna og stengt dører. Ein identisk stol
laga som ein isolert prototype har ikkje hatt denne effekten. Dei
er på same punkt i formrommet, men dei \emph{framtidene} dei møter
er ulike.


\section{Spreiingshypotesen}

\begin{satmark}[Spreiingsrelasjonen]\label{thm:L.8.4}
La $\sigma^2[\gamma]$ vere variansen i formfordelinga rundt ein
gjeven attraktor i $\MC$, og la $\norm{\bm{s}}$ vere den samla
styrken til seleksjonstrykka ved den attraktoren. Då gjeld:
\[
  \sigma^2 \;\propto\; \frac{1}{\norm{\bm{s}}}.
\]
Sterkare seleksjon produserer smalare fordelingar (tettare
stilar); svakare seleksjon produserer breiare fordelingar (meir
variasjon).
\end{satmark}

\begin{proof}[Bevis, steg for steg]
Vi modellerer den stasjonære fordelinga av former nær eit
lokalt maksimum $\bm{x}^*$ som ei Boltzmann-liknande fordeling
(ein modelleringsantaknad, ikkje utleidd frå postulata):
\[
  P(\bm{x}) \;\propto\;
  \exp\!\bigl(\beta \, f(\bm{x})\bigr)
\]
der $\beta > 0$ kontrollerer seleksjonsintensiteten (analog med
invers temperatur i statistisk mekanikk).

\paragraph{Steg 1: Taylor-utvikling rundt maksimumet.}
Nær $\bm{x}^*$ kan vi utvide $f$ til andre orden:
\[
  f(\bm{x}) \approx f(\bm{x}^*)
  - \tfrac{1}{2}
  (\bm{x} - \bm{x}^*)^\top H (\bm{x} - \bm{x}^*)
\]
der $H = -\nabla^2 f(\bm{x}^*)$ er positiv definitt
(sidan $\bm{x}^*$ er eit maksimum).

\paragraph{Steg 2: Sett inn i fordelinga.}
\[
  P(\bm{x}) \propto \exp\!\left(
    -\tfrac{\beta}{2}
    (\bm{x} - \bm{x}^*)^\top H
    (\bm{x} - \bm{x}^*)
  \right).
\]
Dette er ein multivariat normalfordeling med kovariansmatrise
$\Sigma = (\beta H)^{-1}$.

\paragraph{Steg 3: Les av variansen.}
Variansen langs ein vilkårleg akse er
$\sigma_i^2 = (\beta \, \lambda_i)^{-1}$, der $\lambda_i$ er
eigenverdiane til $H$. Sidan $H$ reflekterer kurvaturen til dei
aggregerte seleksjonstrykka, er $\lambda_i \propto \norm{\bm{s}}$,
og vi får:
\[
  \sigma^2 \;\propto\; \frac{1}{\norm{\bm{s}}}. \qedhere
\]
\end{proof}

\paragraph{Intuisjon.}
Tenk på det slik: dersom seleksjonen er sterk --- alle krev same
ting, og det finst lite rom for avvik --- vert formene samla tett
rundt det optimale punktet. Dersom seleksjonen er svak --- ingen
bryr seg særleg --- kan formene spreie seg fritt.


\section{Designhistorisk døme: Kublers formtid}

George Kubler hevda at alt som vert laga er anten ein kopi eller
ein variant av noko som vart laga litt tidlegare. Kvar artefakt er
eit ledd i ei kjede --- ein ``formsekvens'' --- og verdien av ein
artefakt avheng av posisjonen hans i sekvensen, ikkje berre av
eigenskapane hans.

I vårt rammeverk er Kublers observasjon ein konsekvens av
stiavhengig tilpassing. Originalen og kopien okkuperer same punkt
i $\MC$, men dei har ulike stig-funksjonalar $S[\gamma]$:
originalen var det fyrste punktet i ein ny region av formrommet,
medan kopien følgde ein allereie opptrekt rute. Originalen
\emph{opna} rommet; kopien \emph{fylte} det.

Meir presist: originalens $\gamma$ kryssar frå ein attraktor til
ein annan (ein topologisk endring i det personlege
tilpassingslandskapet), medan kopiens $\gamma$ er ein kort,
monoton stig innanfor ein allereie etablert attraktor. Dei to
stigane har ulik ``lengd'' i funksjonalet $S$, og dermed ulik
tilpassingsverdi --- sjølv om dei endar på same stad.


\section{Gjennomarbeidd døme: Spreiingsrelasjonen i praksis}

Vi brukar spreiingsrelasjonen $\sigma^2 \propto 1/\norm{\bm{s}}$
på to kategoriar:

\paragraph{Kategori A: Militærutstyr (sterkt seleksjonstrykk).}
Ein militærhjelm må oppfylle strenge krav til ballistisk vern,
vekt, komfort, og kompatibilitet med anna utstyr. Seleksjonstrykka
er sterke og presist spesifiserte: avvik straffast hardt (i verste
fall fatalt). Vi estimerer ein høg samla seleksjonsstyrke
$\norm{\bm{s}}_A$.

\medskip
Prediksjonen: variansen i hjelmgeometri innanfor ein gjeven
standard (t.d.\ NATO STANAG) skal vere \emph{låg}.

\medskip
Observasjonen: militærhjelmar frå ulike produsentar er
påfallande like. Kompaktheit, kurvatur, massefordeling varierer
minimalt innanfor same standard. Spreiinga er smal.

\paragraph{Kategori B: Dekorative møblar (svakt seleksjonstrykk).}
Ein dekorativ lysestake har svake funksjonelle krav (halde eit lys
oppreist) og sterke estetiske preferansar som varierer mellom
kulturar, epokar og individ. Seleksjonstrykka er svake og
motstridande: det finst ikkje éin ``rett'' form. Vi estimerer ein
låg samla seleksjonsstyrke $\norm{\bm{s}}_B \ll \norm{\bm{s}}_A$.

\medskip
Prediksjonen: variansen i lysestakegeometri skal vere
\emph{høg}.

\medskip
Observasjonen: lysestakar varierer enormt --- frå minimalistiske
sylindrar til barokke kandelabrar. Spreiinga er brei.

\paragraph{Kvantitativt estimat.}
Anta at den effektive seleksjonsintensiteten for militærhjelmar er
$\beta \norm{\bm{s}}_A = 10$ (vilkårleg einingar) og for
dekorative lysestakar $\beta \norm{\bm{s}}_B = 2$. Då:
\[
  \frac{\sigma^2_B}{\sigma^2_A}
  \;=\; \frac{\norm{\bm{s}}_A}{\norm{\bm{s}}_B}
  \;=\; \frac{10}{2} \;=\; 5.
\]
Prediksjonen er at lysestakar har fem gonger høgare geometrisk
varians enn militærhjelmar. Denne prediksjonen er empirisk
testbar: mål kompaktheit, symmetri og kurvatur for eit utval
av kvar kategori og samanlikn variansane.


\section{Øvingar}

\begin{exercise}[Nivå 1]\label{ex:7.1}
Forklar med eigne ord skilnaden mellom å optimere ein
\emph{posisjon} og å optimere ein \emph{stig}. Gje eit døme frå
handverket der prosessen er like viktig som resultatet.
\end{exercise}

\begin{exercise}[Nivå 1]\label{ex:7.2}
Ein original Wegner-stol og ein fabrikkprodusert replika har
identiske mål. Bruk omgrepet stiavhengig tilpassingsverdi til å
forklare kvifor dei har ulik status i formgjevingsøkologien.
\end{exercise}

\begin{exercise}[Nivå 2]\label{ex:7.3}
I Euler--Lagrange-likninga for designtrajektoriar:
$\frac{\partial g}{\partial \gamma}
- \frac{d}{dt}\frac{\partial g}{\partial \dot{\gamma}} = 0$,
kva svarar kvart av dei to ledda til i ein konkret designprosess?
Gje eit døme der det fyrste leddet dominerer og eit der det andre
dominerer.
\end{exercise}

\begin{exercise}[Nivå 2]\label{ex:7.4}
Spreiingshypotesen seier at $\sigma^2 \propto 1/\norm{\bm{s}}$.
Vel to funksjonelle klassar --- éin med sterkt seleksjonstrykk og
éin med svakt --- og prediker forholdet mellom variansane deira.
Korleis ville du teste prediksjonen empirisk?
\end{exercise}

\begin{exercise}[Nivå 3]\label{ex:7.5}
Teoremet om at stigen modifiserer landskapet (Teorem~\ref{thm:L.8.3})
har ein konkret konsekvens: to designtradisjonar som startar frå
same punkt men tek ulike ruter, kan ende opp med heilt ulike sett
av tilgjengelege optima. Konstruer eit tenkt døme: spesifiser eit
formrom med to dimensjonar, to tradisjonar med ulike starttrajektoriar,
og vis korleis landskapsminnefunksjonen $\eta_i$ produserer ulike
framtidige landskap.
\end{exercise}

\begin{exercise}[Nivå 3]\label{ex:7.6}
Kublers formsekvens-omgrep svarar til stigen $\gamma$ i
formrommet. Men Kubler antok at sekvensane er \emph{lineære}
kjeder. I vårt rammeverk er stigen ein kontinuerleg kurve i eit
$n$-dimensjonalt rom, som kan forgreine seg. Drøft kva
konsekvensar det har for Kublers typologi at formrommet er
fleirdimensjonalt, ikkje éindimensjonalt.
\end{exercise}


%% =====================================================================
%%  KAPITTEL 8: EMPIRISK --- MATERIALSIGNATURAR
%% =====================================================================
\chapter{Empirisk: Materialsignaturar}
\label{chap:empirisk}

\begin{quote}
\itshape
Materialet er ikkje nøytralt. Det kanaliserer forma.
\end{quote}

\section{Frå teori til data}

Kapitla~1--7 bygde opp eit deduktivt rammeverk: frå postulat til
teorem, frå definisjonar til konsekvensar. Alt kviler på
antaknadar som vi \emph{hevdar} er sanne, men som vi endå ikkje
har testa mot data.

Dette kapittelet gjer noko annleis. Det presenterer eit empirisk
resultat --- det sterkaste enkelttestresultatet av rammeverket ---
og viser korleis det passar inn. Kapittelet følgjer ikkje logisk
frå dei føregåande proposisjonane; det \emph{illustrerer} dei.


\section{Postulat om geometrisk signatur}

\begin{postmark}[Geometrisk materialsignatur]\label{post:L.6.1}
For kvart materiale $m$ brukt i ein funksjonell klasse $C$ er den
betinga fordelinga av geometriske eigenskapar gjeve materiale
ikkje-uniform:
\[
  P(\bm{x} \mid m) \;\neq\; \mathrm{Uniform}(\MC)
  \qquad \forall\; m.
\]
Kvart materiale induserer ein karakteristisk geometrisk
\emph{profil}: ein region i formrommet som det favoriserer.
\end{postmark}

\paragraph{Falsifiseringsvilkår.}
Postulatet fell om same geometriske distribusjon kan observerast
uavhengig av materiale, innanfor same funksjonelle klasse og
same sett av andre seleksjonstrykk.

\paragraph{Kva betyr dette?}
Påstanden er at materialet \emph{kanaliserer} forma. Mahogni
trekkjer mot andre geometriar enn stål, ikkje fordi designaren
vel det, men fordi dei innbygde eigenskapane til materialet ---
brotstyrke, bøyelighet, bearbeidingskarakter --- favoriserer visse
konfigurasjonar og motset seg andre. Denne kanaliseringa er
målbar og systematisk.


\section{Informasjonsteoremet}

\begin{satmark}[Informasjonsteoremet]\label{thm:L.6.2}
Innanfor ein funksjonell klasse $C$ med minst to materialkategoriar
ber materialet strengt meir gjensidig informasjon med realisert
geometri enn funksjonen gjer:
\[
  \MI(G;\, M) \;>\; \MI(G;\, \mathrm{Func}) \;=\; 0.
\]
\end{satmark}

\begin{proof}[Bevis, steg for steg]

\paragraph{Del~1: $\MI(G;\, \mathrm{Func}) = 0$.}%
\begin{enumerate}
  \item \textbf{Premiss.}\;
    Innanfor ein funksjonell klasse $C$ deler alle objekt same
    funksjon per definisjon (kap.~2).
  \item \textbf{Konsekvens.}\;
    Variabelen $\mathrm{Func}$ er \emph{konstant} innanfor $C$.
    Ein konstant variabel har null entropi:
    $\Ent(\mathrm{Func}) = 0$.
  \item \textbf{Gjensidig informasjon.}\;
    Gjensidig informasjon er definert som
    $\MI(G;\, \mathrm{Func}) = \Ent(G) -
    \Ent(G \mid \mathrm{Func})$.
    Sidan $\mathrm{Func}$ er konstant, gjev betinginga ingen ny
    informasjon:
    $\Ent(G \mid \mathrm{Func}) = \Ent(G)$.
  \item \textbf{Resultat.}\;
    $\MI(G;\, \mathrm{Func}) = \Ent(G) - \Ent(G) = 0$.
\end{enumerate}

\paragraph{Del~2: $\MI(G;\, M) > 0$.}%
\begin{enumerate}
  \item \textbf{Premiss.}\;
    Frå postulat~\ref{post:L.6.1}:
    $P(\bm{x} \mid m) \neq P(\bm{x})$ for minst eitt
    materiale~$m$. Det vil seie at den materiale-betinga
    fordelinga skil seg frå marginalfordelinga.
  \item \textbf{Konsekvens.}\;
    Dersom $P(\bm{x} \mid m) \neq P(\bm{x})$, er $G$ og $M$
    \emph{ikkje} uavhengige: $G \not\perp\!\!\!\perp M$.
  \item \textbf{Fundamental eigenskap.}\;
    Gjensidig informasjon er null om og berre om variablane er
    uavhengige:
    $\MI(G;\, M) = 0 \;\iff\; G \perp\!\!\!\perp M$.
  \item \textbf{Resultat.}\;
    Sidan $G \not\perp\!\!\!\perp M$, har vi $\MI(G;\, M) > 0$.
\end{enumerate}

\paragraph{Kombinert:}
$\MI(G;\, M) > 0 = \MI(G;\, \mathrm{Func})$. Materialet ber meir
informasjon om forma enn funksjonen gjer. \qed
\end{proof}

\paragraph{Kvifor er dette viktig?}
Resultatet er eit direkte motargument mot det deterministiske
aksiomet ``form følgjer funksjon''. Innanfor ein funksjonell
klasse --- der funksjonen er \emph{identisk} for alle objekt ---
forklarer funksjonen \emph{null} av den geometriske variasjonen.
Materialet forklarer ein målbar, positiv del. Funksjonen
\emph{avgrensar}; materialet \emph{formar}.


\section{STOLAR-datasettet}

Beviset ovanfor er abstrakt. Tala under er konkrete. Dei kjem frå
STOLAR-datasettet: 93~stolar frå eit norsk nasjonalmuseum,
daterte frå 1280 til 2024, fordelte på fire materialkategoriar.

\subsection{Geometriske profilar per materiale}

For kvar stol i datasettet er tre geometriske eigenskapar målte:
\emph{kompaktheit} (kor tett forma er pakka), \emph{symmetri}
(grad av bilateralt spegelsymmetri), og \emph{kurvatur}
(gjennomsnittleg overflatebuktning). Resultata per materiale:

\medskip
\begin{center}
\begin{tabular}{@{}lccc@{}}
\toprule
\textsc{Materiale} & \textsc{Kompaktheit}
  & \textsc{Symmetri} & \textsc{Kurvatur} \\
\midrule
Mahogni    & 0{,}381 & 0{,}129 & 0{,}105 \\
Stål       & 0{,}437 & 0{,}118 & 0{,}075 \\
Kryssfiner & 0{,}455 & 0{,}104 & 0{,}090 \\
Bøk        & 0{,}445 & 0{,}101 & 0{,}101 \\
\bottomrule
\end{tabular}
\end{center}

\medskip
Mønstra er tydelege:
\begin{itemize}[nosep]
  \item \textbf{Mahogni er minst kompakt.} Materialet si høge
    brotstyrke let formgjevaren spreie massen i slankare lemer.
  \item \textbf{Stål har lågast kurvatur.} Metallet tenderer mot
    det lineære --- bøying krev spesialiserte prosessar.
  \item \textbf{Kryssfiner er mest kompakt.} Laminatet inviterer
    til skalformer og lukka profilar.
  \item \textbf{Bøk har høgast kurvatur.} Dampbøying opnar for
    rundare former.
\end{itemize}

Desse profilane er ikkje styrte av funksjonen --- alle 93 objekta
er stolar. Det er materialets innbygde eigenskapar som kanaliserer
forma. Effektstorleikane er store: ull/massefordeling $d = 1{,}60$;
metall/høgde $d = 1{,}47$; skumplast/massefordeling $d = 1{,}40$
(Cohen's $d$). Distribusjonane overlappar knapt.

\subsection{Empiriske informasjonsverdiar}

Frå datasettet:
\begin{align*}
  \MI(G;\, M) &= 0{,}382 \;\text{bit}, \\
  \MI(G;\, \mathrm{Tid}) &= 0{,}168 \;\text{bit}, \\
  \MI(G;\, \mathrm{Geografi}) &= 0{,}079 \;\text{bit}, \\
  \MI(G;\, \mathrm{Func}) &= 0 \;\text{bit (per definisjon).}
\end{align*}
Materialet er den sterkaste prediktoren. Det ber meir enn dobbelt
so mykje informasjon om forma som tida gjer, og nesten fem gonger
so mykje som geografien.

\subsection{Materialentropien over tid}

I perioden 1825--1850 utgjorde mahogni 93~prosent av alle
registrerte norskproduserte stolar i nasjonalmuseet. Material\-entropien --- eit informasjonsteoretisk mål på mangfald --- kollapsa
frå 2{,}5 til 1{,}5~bit. Heile designlandskapet konvergerte mot
éin topp: eit klassisk lokalt optimum der eitt seleksjonstrykk var
so dominant at alt anna forsvann.

Denne ``mahognitoppen'' illustrerer spreiingsrelasjonen
(Teorem~\ref{thm:L.8.4}) i praksis: ekstremt sterkt
seleksjonstrykk --- materialtilgang, prestisje, kolonialt
handelsnettverk --- produserte ein ekstremt smal distribusjon.


\section{Designhistorisk døme: Mahognitoppen}

Mellom 1750 og 1850 dominerte mahogni den norske møbelproduksjonen
i ein grad som er vanskeleg å overdrive. Materialet kom frå
Karibia og Mellom-Amerika via eit globalt handelsnettverk forma av
kolonialisme, sjøfart og europeisk prestisjeøkonomi. Kvar stol
laga i norsk mahogni er eit komprimert verdskart: trevyrket
knyter Christiania til Kingston, snikkaren til plantasjearbeidaren,
den lokale smaken til den globale materialstraumen.

42~prosent av mahognistolane i det norskproduserte delsettet
kombinerer importert mahogni med lokalt tre: mahogni utanpå, furu
inni. Stolen er \emph{kolonialt finert}: globalt prestisjetømmer
dekkjer lokal struktur. Materialkurva er eit komprimert verdskart.
Kvar stol er ein geopolitisk artefakt.

Rammeverket fangar denne dynamikken presist: materialet definerer
ein geometrisk profil (låg kompaktheit, middels symmetri, middels
kurvatur), og det kulturelle seleksjonstrykket (prestisje, mote,
tilgang) konsentrerer populasjonen rundt denne profilen til
variansen kollapsar. Når mahognien forsvinn --- fyrst grunna
overhogging, seinare grunna stilskifte --- opnar formrommet seg att,
og entropien stig.


\section{Gjennomarbeidd døme: Rekn ut gjensidig informasjon}

Vi reknar ut gjensidig informasjon $\MI(G;\, M)$ for eit forenkla
datasett med 2~materialkategoriar og 2~geometriske kategoriar.

\paragraph{Datasettet.}
20~objekt i funksjonell klasse ``stol'', fordelte slik:

\medskip
\begin{center}
\begin{tabular}{@{}lcc|c@{}}
\toprule
& \textsc{Kompakt} & \textsc{Open} & \textsc{Sum} \\
\midrule
Tre & 8 & 2 & 10 \\
Stål & 3 & 7 & 10 \\
\midrule
Sum & 11 & 9 & 20 \\
\bottomrule
\end{tabular}
\end{center}

\paragraph{Steg 1: Rekn ut marginalfordeling.}
\begin{align*}
  P(\text{Tre}) &= 10/20 = 0{,}5, &
  P(\text{Stål}) &= 10/20 = 0{,}5, \\
  P(\text{Kompakt}) &= 11/20 = 0{,}55, &
  P(\text{Open}) &= 9/20 = 0{,}45.
\end{align*}

\paragraph{Steg 2: Rekn ut fellesfordeling.}
\begin{align*}
  P(\text{Tre, Kompakt}) &= 8/20 = 0{,}40, \\
  P(\text{Tre, Open}) &= 2/20 = 0{,}10, \\
  P(\text{Stål, Kompakt}) &= 3/20 = 0{,}15, \\
  P(\text{Stål, Open}) &= 7/20 = 0{,}35.
\end{align*}

\paragraph{Steg 3: Rekn ut gjensidig informasjon.}
Formelen er:
\[
  \MI(G;\, M) = \sum_{g, m} P(g, m) \,
  \log_2 \frac{P(g, m)}{P(g) \, P(m)}.
\]

Vi reknar ut kvart ledd:
\begin{align*}
  &P(\text{Tre, K}) \log_2
    \frac{P(\text{Tre, K})}{P(\text{K})\,P(\text{Tre})}
  = 0{,}40 \cdot \log_2 \frac{0{,}40}{0{,}55 \cdot 0{,}50}
  = 0{,}40 \cdot \log_2 1{,}4545 \\
  &\quad= 0{,}40 \cdot 0{,}5406 = 0{,}2162, \\[4pt]
  %
  &P(\text{Tre, O}) \log_2
    \frac{P(\text{Tre, O})}{P(\text{O})\,P(\text{Tre})}
  = 0{,}10 \cdot \log_2 \frac{0{,}10}{0{,}45 \cdot 0{,}50}
  = 0{,}10 \cdot \log_2 0{,}4444 \\
  &\quad= 0{,}10 \cdot (-1{,}1699) = -0{,}1170, \\[4pt]
  %
  &P(\text{Stål, K}) \log_2
    \frac{P(\text{Stål, K})}{P(\text{K})\,P(\text{Stål})}
  = 0{,}15 \cdot \log_2 \frac{0{,}15}{0{,}55 \cdot 0{,}50}
  = 0{,}15 \cdot \log_2 0{,}5454 \\
  &\quad= 0{,}15 \cdot (-0{,}8745) = -0{,}1312, \\[4pt]
  %
  &P(\text{Stål, O}) \log_2
    \frac{P(\text{Stål, O})}{P(\text{O})\,P(\text{Stål})}
  = 0{,}35 \cdot \log_2 \frac{0{,}35}{0{,}45 \cdot 0{,}50}
  = 0{,}35 \cdot \log_2 1{,}5556 \\
  &\quad= 0{,}35 \cdot 0{,}6376 = 0{,}2232.
\end{align*}

\paragraph{Steg 4: Summer.}
\[
  \MI(G;\, M) = 0{,}2162 - 0{,}1170 - 0{,}1312 + 0{,}2232
  = 0{,}1912 \;\text{bit}.
\]

\paragraph{Tolking.}
Resultatet $\MI(G;\, M) = 0{,}19$~bit er strengt positivt: å vite
materialet gjev informasjon om geometrien. Tresstolar tenderer mot
kompakte former; stålstolar mot opne. Funksjonen (``stol'') er
konstant og gjev $\MI(G;\, \mathrm{Func}) = 0$. Materialet
forklarer variasjonen som funksjonen ikkje kan.


\section{Øvingar}

\begin{exercise}[Nivå 1]\label{ex:8.1}
Kvifor er $\MI(G;\, \mathrm{Func}) = 0$ innanfor ein funksjonell
klasse? Forklar med eigne ord, utan formlar.
\end{exercise}

\begin{exercise}[Nivå 1]\label{ex:8.2}
Sjå på dei geometriske profilane per materiale i
STOLAR-datasettet. Kva profil ville du forvente for \emph{plast}?
Grunngjev svaret med referanse til plastens fysiske eigenskapar.
\end{exercise}

\begin{exercise}[Nivå 2]\label{ex:8.3}
Utvid det gjennomarbeidde dømet: legg til ein tredje material\-kategori
(``plast'') med 10~objekt, der 5~er kompakte og 5~er opne. Rekn
ut den nye $\MI(G;\, M)$ for 3~material og 2~geometriar. Kva
skjer med verdien? Kvifor?
\end{exercise}

\begin{exercise}[Nivå 2]\label{ex:8.4}
Materialentropien over tid viser ein kollaps under mahognitoppen.
Formuler dette som eit tilfelle av spreiingsrelasjonen
(Teorem~\ref{thm:L.8.4}): kva er $\norm{\bm{s}}$, og kvifor
veks det i perioden 1750--1850?
\end{exercise}

\begin{exercise}[Nivå 3]\label{ex:8.5}
Designa eit empirisk eksperiment for å teste postulat~\ref{post:L.6.1}
på ein annan funksjonell klasse enn stolar (t.d.\ knivar, koppar,
eller syklar). Spesifiser: (a)~dei geometriske eigenskapane du
ville måle, (b)~materialkategoriane, (c)~korleis du ville
kontrollere for andre seleksjonstrykk, og (d)~kva statistisk test
du ville bruke.
\end{exercise}

\begin{exercise}[Nivå 3]\label{ex:8.6}
Informasjonsteoremet (Teorem~\ref{thm:L.6.2}) kviler på postulat~%
\ref{post:L.6.1}. Konstruer eit hypotetisk motdøme: ein
funksjonell klasse og eit sett av materialkategoriar der
$P(\bm{x} \mid m) = P(\bm{x})$ for alle $m$. Kva ville det krevje
i praksis? Kvifor er det usannsynleg?
\end{exercise}


%% =====================================================================
%%  KAPITTEL 9: KONSEKVENSAR OG FRAMTID
%% =====================================================================
\chapter{Konsekvensar og framtid}
\label{chap:konsekvensar}

\begin{quote}
\itshape
Ingen form er endeleg. Kvart objekt er eit augeblinksbilete
av eit dynamisk system.
\end{quote}


\section{Provisorisk kompromiss}

\begin{satmark}[Provisorisk kompromiss]\label{thm:L.9.1}
Kvar realisert form er eit \emph{provisorisk} kompromiss:
optimalt for seleksjonstrykka på tidspunktet det vart skapt,
men generisk suboptimalt for alle seinare konfigurasjonar av trykk.
\end{satmark}

\begin{proof}[Grunngjeving]
Frå kompromissteoremet (kap.~2) veit vi at kvar realisert form
$\bm{x} \in B$ balanserer motstridande seleksjonstrykk utan å
maksimere alle samstundes. Frå postulatet om dynamisk landskap
(kap.~4) veit vi at trykka $s_i(\bm{x}, t)$ endrar seg over tid.

La $\bm{x}^*(\tau)$ vere det lokale maksimumet av $f(\cdot, \tau)$
nærmast $\bm{x}$ på skapingstidspunktet~$\tau$. På eit seinare
tidspunkt $t > \tau$ har trykka skifta:
$\bm{s}(\cdot, t) \neq \bm{s}(\cdot, \tau)$.
Punktet $\bm{x}^*(\tau)$ er generisk \emph{ikkje} eit kritisk
punkt av $f(\cdot, t)$, sidan:
\[
  \grad f(\bm{x}^*(\tau),\, t) =
  \int_\tau^t \frac{\partial}{\partial t'}
  \grad f(\bm{x}^*(\tau),\, t')\, dt'
  \neq \bm{0}
\]
med mindre integralet er identisk null (eit ikkje-generisk
vilkår). Forma som var optimal ved $\tau$ er ikkje lenger optimal
ved $t$. Ingen form er endeleg; kvart objekt er eit augeblinksbilete
av eit dynamisk system.
\end{proof}

\paragraph{Konsekvens.}
Kvar stol, kvart hus, kvar organisme er midlertidig. Ikkje fordi
det er dårleg laga, men fordi verda rundt det endrar seg. Ein
``tidlaus'' design er ein design der seleksjonstrykka tilfeldigvis
har endra seg lite --- ikkje ein design som har transcendert
endring.


\section{Varige tradisjonar}

\begin{satmark}[Varige tradisjonar]\label{thm:L.9.2}
Tradisjonar med fleire uavhengige adaptive mekanismar (fleire
seleksjonstrykk adresserte av kompetente sub-navigatorar) er meir
robuste mot landskapsendringar.
\end{satmark}

\begin{proof}[Grunngjeving]
La ein tradisjon $\mathcal{T}$ vere ein populasjon av navigatorar
som adresserer $k_\mathcal{T}$ ulike seleksjonstrykk gjennom
uavhengige adaptive mekanismar. Når landskapet skiftar (kap.~4),
endrar typisk berre eit delsett av trykka seg samstundes.

Ein tradisjon med $k_\mathcal{T}$ uavhengige mekanismar kan
respondere på eit skifte i trykk $s_i$ utan å forstyrre responsen
på trykk $s_j$ ($j \neq i$), fordi mekanismane er uavhengige.

Samanlikn tradisjonar $\mathcal{T}_1$ med $k_1$ mekanismar og
$\mathcal{T}_2$ med $k_2 > k_1$ mekanismar. Under eit
landskapsskifte som påverkar eitt trykk, justerer $\mathcal{T}_2$
berre $\Delta k / k_2 < \Delta k / k_1$ av den totale adaptive
kapasiteten sin. Det større repertoaret gjev redundans.
\end{proof}

\paragraph{I praksis.}
Ei handverkstradisjon som kombinerer materialforståing, kulturell
sensitivitet og individuell innovasjon tilpassar seg betre enn
ei som kviler på éin einaste mekanisme. Ei masseproduksjonslinje
som berre responderer på kostnadstrykk er sårbar når smaken
endrar seg. Ein AI-generator som berre optimerer for popularitet
kollapsar når brukarane skiftar preferansar.


\section{Falsifiserbarheit}

\begin{satmark}[Rammeverket kan felle seg sjølv]%
\label{thm:L.9.3}
Rammeverket i denne boka er falsifiserbart: dersom eitt einaste
seleksjonstrykk $s_1$ er tilstrekkeleg til å forklare all
observert formvariasjon innanfor ein funksjonell klasse $C$, fell
postulat~2.2 (fleire trykk), og med det fell teorema~2.21, 2.31,
2.32, 3.3, 3.41, og~\ref{thm:L.9.1}.
\end{satmark}

\begin{proof}[Den deduktive kjeda]
\begin{enumerate}[nosep]
  \item Postulat om fleire trykk $\implies$ Teorem om variasjon
    under konstant funksjon.
  \item Postulat om fleire trykk + postulat om motstridande
    gradientar $\implies$ kompromissteoremet $\implies$ teorem
    om fleire kompromiss.
  \item Kompromissteoremet + fleire kompromiss $\implies$ teorem
    om fleire toppar $\implies$ stildefinisjonen.
  \item Fleire toppar + stildefinisjonen $\implies$ korollar om
    analytikaravhengigheit.
  \item Dynamisk landskap + kompromissteoremet $\implies$
    provisorisk kompromiss (Teorem~\ref{thm:L.9.1}).
\end{enumerate}
Dersom nokon av dei sju postulata kan motbevisast empirisk,
kollapsar den tilsvarande greina av det deduktive treet.
Rammeverket er gyldig \emph{nettopp fordi} det kan feile.

Det deterministiske aksiomet ``form følgjer funksjon'' var
ufalsifiserbart og dermed ikkje ein vitskapeleg påstand. Denne
teorien gjer kvantitative prediksjonar
(Teorem~\ref{thm:L.8.4}, Postulat~\ref{post:L.6.1}) som kan
testast mot data. At teksten \emph{kan} fellast er grunnen til at
ho er gyldig.
\end{proof}


\section{Refleksiv navigasjon}

Alle navigatorane vi har diskutert til no --- termostaten,
algoritmen, evolusjonen --- har éin ting til felles: dei kan ikkje
endre sine eigne mål. Termostaten følgjer settpunktet.
Gradient-descent følgjer tapsfunksjonen. Evolusjonen følgjer
tilpassingslandskapet. Ingen av dei kan \emph{spørje seg sjølv}
kvifor dei favoriserer ei viss form, og endre kursen fordi svaret
ikkje held mål.

Menneskelege designarar kan det.

Ein handverkar kan spørje seg kvifor ho favoriserer ei viss
stolform, og endre kursen fordi ein refleksjon over eigne mål har
endra kva som tel som ``betre''. Ein diffusjonsmodell kan ikkje
det. Refleksiviteten betyr ikkje at navigatoren opphevar
gradientane --- ho er framleis i landskapet, framleis underlagt
trykk. Men ho kan \emph{velje} å gå mot gradienten: å akseptere
lågare tilpassing fordi ei refleksjon over eigne mål har endra
kva som tel.

Klasse~D-navigatorar (dei som responderer på grunnar, ikkje berre
på signal) er ikkje berre raskare eller betre informerte
gradientfølgjarar. Dei er system som kan \emph{omskrive sin eigen
tilpassingsfunksjon} --- ikkje berre sine eigne kriterium for
kva det vil seie å nå dei.

Denne kapasiteten er det som gjer designutdanning mogleg og
naudsynt: studentane lærer ikkje berre å navigere betre, men å
\emph{endre kva navigasjon betyr}.


\section{Fire praksisar analyserte}

Rammeverket er ikkje berre abstrakt teori. Det er eit verktøy for
å forstå konkrete formgjevingspraksisar. Vi analyserer fire:

\subsection{Handverksdesign}

\paragraph{Navigatoren.}
Handverkar + verktøy + materiale. Eit trippel der den taktile
tilbakemeldinga frå materialet er ein integrert del av
styringsvektoren~$\alpha$.

\paragraph{Formrommet.}
Rommet av former som verktøyet og materialet kan realisere. Ein
snikkar i bøk har eit anna formrom enn ein keramikar i steingods,
sjølv om begge kan lage skåler.

\paragraph{Seleksjonstrykka.}
Oppdragsgjevar, tradisjon, materialøkonomi, ergonomi, marknad,
kultur. Fleire uavhengige trykk, ofte i konflikt.

\paragraph{Tilpassingslandskapet.}
Komplekst, multimodalt, sterkt forma av lokalkultur. Toppane
svarar til etablerte stilar; dalane til former ingen bryr seg om.

\paragraph{Eksplorasjon og eksploitering.}
Eksplorasjon = eksperiment, feilgrep. Eksploitering =
meisterarbeid, gjentaking, perfeksjonering. Handverkaren
balanserer mellom å prøve noko nytt og å meistre det kjende.

\paragraph{Styrke og svakheit.}
Styrken er materialforståing og refleksiv kapasitet. Svakheita er
dekning: handverkaren utforskar berre ein liten del av
formrommet.

\subsection{Industridesign}

\paragraph{Navigatoren.}
Designar + ingeniør + produksjonslinje. Navigasjonen er
distribuert: ingen einskild person kontrollerer alle
dimensjonane.

\paragraph{Formrommet.}
Rommet av former som kan masseproduseras med gjeven verktøykostnad.
Sterkare avgrensingar enn handverket, men større volum.

\paragraph{Seleksjonstrykka.}
Einheitskostnad, ergonomi, estetikk, produksjonseffektivitet,
logistikk, regulering. Kostnads- og produksjonstrykka dominerer
ofte.

\paragraph{Eksplorasjon og eksploitering.}
Eksplorasjon = prototypar, brukartestar. Eksploitering =
optimering av produksjonslinja, standardisering.

\paragraph{Modulering.}
Industridesign brukar \emph{modular dekomponering}: systemet vert
delt i komponentar som kan varierast uavhengig. Kvar modul er ein
sub-navigator med eige formrom. Emergent navigasjon oppstår
mellom modulane.

\subsection{Parametrisk og algoritmisk design}

\paragraph{Navigatoren.}
Algoritme + designar (som definerer parameterrommet og
tilpassingsfunksjonen). Designaren navigerer ikkje direkte i
formrommet, men i \emph{metarommet} av parametrar.

\paragraph{Formrommet.}
Det parametrisk definerte delrommet av $\MC$. Mykje smalare enn
det fulle formrommet, men eksplisitt og systematisk utforska.

\paragraph{Seleksjonstrykka.}
Eksplisitt formulerte tilpassingsfunksjonar. Styrken er at
trykka er \emph{spesifiserte}, ikkje implisitte. Svakheita er at
det som ikkje kan spesifiserast, ikkje vert optimert.

\paragraph{Eksplorasjon og eksploitering.}
Eksplorasjon = parametervariasjon, Monte Carlo-sampling,
genetiske algoritmar. Eksploitering = gradient descent,
konvergering mot eit lokalt optimum.

\paragraph{Latent rom.}
I parametrisk design er ``latent rom'' metarommet: det rommet av
alle parametervektorar. Gradient descent i dette rommet svarar
til rørsle i formrommet. Kvar parametervektor $\bm{p}$ definerer
ein posisjon $\bm{x}(\bm{p}) \in \MC$ gjennom ein avbildingsfunksjon.

\subsection{Generativ AI-design}

\paragraph{Navigatoren.}
Trena nevralt nettverk + menneskeleg evaluering. Nettverket har
lært ei implisitt framstilling av formrommet frå store datasett.

\paragraph{Formrommet.}
Latent rom (latent space): det indre representasjonsrommet til
nettverket. Ikkje direkte tolkbart, men navigerbart gjennom
interpolasjon og sampling.

\paragraph{Seleksjonstrykka.}
Treningsdata (implisitt) + prompts (eksplisitt) + menneskeleg
evaluering. Tilpassingslandskapet er den implisitte fordelinga av
``gode'' former i treningssettet.

\paragraph{Temperatur og spreiing.}
Spreiingsrelasjonen (Teorem~\ref{thm:L.8.4}) har ein direkte
parallell i generativ AI: \emph{temperaturparameteren}. Høg
temperatur = meir eksplorasjon, breiare spreiing, meir variasjon.
Låg temperatur = meir eksploitering, smalare spreiing,
deterministisk utdata. Seleksjonsintensiteten $\beta$ i
spreiingsrelasjonen svarar direkte til invers temperatur
$1/T$ i samplingsalgoritmen.

\paragraph{Diffusjonsmodellar.}
Diffusjonsmodellar navigerer gjennom \emph{denoising}: dei
startar frå støy og konvergerer trinnvis mot ein lært attraktor.
Kvart steg er eit steg opp i det lærde tilpassingslandskapet.
Prosessen er ein konkret realisering av gradient-ascent i eit
approksimert tilpassingslandskap.

\paragraph{Styrke og svakheit.}
Styrken er dekning: nettverket har sett store delar av
formrommet. Svakheita er tolkbarheit: det kan ikkje forklare
\emph{kvifor} det valde ein gjeven posisjon.

\subsection{Evaluering langs fire aksar}

Kvar formgjevingsprosess kan evaluerast langs fire aksar:

\medskip
\begin{center}
\begin{tabular}{@{}lcccc@{}}
\toprule
& \textsc{Handverk} & \textsc{Industri}
  & \textsc{Parametrisk} & \textsc{Gen. AI} \\
\midrule
Dekning & Låg & Middels & Middels & Høg \\
Effektivitet & Låg & Middels & Høg & Høg \\
Robustheit & Høg & Middels & Låg & Låg \\
Tolkbarheit & Høg & Middels & Middels & Låg \\
\bottomrule
\end{tabular}
\end{center}

\medskip
\begin{description}[nosep]
  \item[Dekning:] Kor stor del av formrommet utforskar substratet?
  \item[Effektivitet:] Kor raskt finn substratet eit lokalt
    optimum?
  \item[Robustheit:] Kor godt fungerer substratet når
    tilpassingslandskapet endrar seg?
  \item[Tolkbarheit:] Kan ein forstå \emph{kvifor} substratet
    valde ein posisjon?
\end{description}

\medskip
Ingen substrat er optimalt langs alle fire aksane. Å velje substrat
er å velje kva ein prioriterer.


\section{Kva betyr dette for designutdanning?}

Dersom rammeverket i denne boka er rett, har det konsekvensar for
korleis vi underviser design:

\begin{enumerate}
  \item \textbf{Formrommet må ekspliserast.}
    Studentar bør lære å identifisere og kartleggje formrommet for
    oppgåva si --- ikkje berre å ``utforske'' intuitivt, men å
    forstå kva dimensjonar som finst og kvar dei er.

  \item \textbf{Seleksjonstrykka må namngjevast.}
    Kvar designavgjerd er ein vekting av motstridande trykk.
    Studentar bør lære å identifisere trykka, ikkje berre å
    ``balansere''.

  \item \textbf{Kompromisset må anerkjennast.}
    Ingen form maksimerer alt samstundes. Å lære design er å lære
    å gjere medvitne kompromiss, ikkje å streve etter ein illusjon
    av perfeksjon.

  \item \textbf{Substratval er verktøyval er epistemologisk val.}
    Å velje å teikne for hand, å bruke parametrisk modellering, å
    promte ein diffusjonsmodell --- kvart val opnar delar av
    formrommet og stengjer andre. Studenten bør forstå \emph{kva}
    ho gjev opp når ho vel eit substrat.

  \item \textbf{Refleksjon er ikkje luksus, men navigasjonsmodus.}
    Refleksiv navigasjon --- evna til å endre eigne mål --- er det
    som skil designaren frå termostaten. Designutdanning som ikkje
    utviklar refleksiv kapasitet, produserer sofistikerte
    termostatar, ikkje designarar.
\end{enumerate}


\section{Øvingar}

\begin{exercise}[Nivå 1]\label{ex:9.1}
Forklar med eigne ord kva ``provisorisk kompromiss'' betyr. Gje eit
døme på ein artefakt som var optimal på eit tidspunkt men er
suboptimal no.
\end{exercise}

\begin{exercise}[Nivå 1]\label{ex:9.2}
Kvifor er falsifiserbarheit ein styrke, ikkje ein svakheit, for eit
teoretisk rammeverk?
\end{exercise}

\begin{exercise}[Nivå 2]\label{ex:9.3}
Vel éin av dei fire praksisane (handverk, industri, parametrisk,
generativ AI) og analyser ei konkret designoppgåve (t.d.\ å
designe ein lampe) ved å spesifisere: (a)~navigatoren, (b)~formrommet,
(c)~seleksjonstrykka, (d)~tilpassingslandskapet, og
(e)~eksplorasjon vs. eksploitering.
\end{exercise}

\begin{exercise}[Nivå 2]\label{ex:9.4}
Teoremet om varige tradisjonar (Teorem~\ref{thm:L.9.2}) seier at
tradisjonar med fleire adaptive mekanismar er meir robuste. Finn
eit døme på ein handverkstradisjon som har overlevd i meir enn
200~år, og identifiser minst tre uavhengige adaptive mekanismar
som har bidrege til overlevinga.
\end{exercise}

\begin{exercise}[Nivå 3]\label{ex:9.5}
Refleksiv navigasjon inneber at navigatoren kan endre
tilpassingsfunksjonen sin. Formuler dette matematisk: kva betyr det
formelt at $\Phi$ er ein funksjon av $\gamma([0, t])$? Kva
konsekvensar får det for Euler--Lagrange-likninga
(Teorem~\ref{thm:L.8.2})?
\end{exercise}

\begin{exercise}[Nivå 3]\label{ex:9.6}
Temperaturparameteren i generativ AI svarar til invers
seleksjonsintensitet $\beta$ i spreiingsrelasjonen. Design eit
eksperiment der du brukar ein biletgenereringsmodell med ulike
temperaturverdiar og måler den geometriske variansen i utdata.
Kva predikerer spreiingsrelasjonen? Korleis ville du teste det?
\end{exercise}

\begin{exercise}[Nivå 3]\label{ex:9.7}
Tabellen over fire praksisar evaluert langs fire aksar er ein
\emph{påstand}, ikkje eit bevis. Vel éin celle i tabellen (t.d.\
``robustheit for generativ AI = låg'') og konstruer eit argument
for eller mot den vurderinga. Kva data treng du for å avgjere
spørsmålet empirisk?
\end{exercise}


%% =====================================================================
%%  APPENDIKS
%% =====================================================================
\appendix

%% =====================================================================
%%  APPENDIKS A: KORRESPONDANSETABELL
%% =====================================================================
\chapter{Korrespondansetabell}
\label{app:korrespondanse}

Tabellen nedanfor knyter kvart formelt element i denne læreboka
til den tilsvarande posisjonen i \emph{Tractatus Formae} og i
\emph{Formlære: Mathematica}. Nummereringa i traktaten brukar
desimalproposisjonar; \textsc{Mathematica} brukar
seksjon.element; læreboka brukar kapittelnummerering.

\medskip
{\small
\begin{longtable}{@{}p{42mm}p{24mm}p{22mm}@{}}
\toprule
\textsc{Lærebok} & \textsc{Tractatus} & \textsc{Mathematica} \\
\midrule
\endfirsthead
\toprule
\textsc{Lærebok} & \textsc{Tractatus} & \textsc{Mathematica} \\
\midrule
\endhead
\bottomrule
\endfoot

Kap.~1: Formrommet & 1.1--1.14 & Def.~1.1--Prop.~1.4 \\
Kap.~2: Seleksjonstrykk & 2.1--2.4 & Def.~2.1--Thm.~2.8 \\
Kap.~3: Tilpassingslandskapet & 3.1--3.6 & Def.~3.1--Def.~3.6 \\
Kap.~4: Dynamikk & 4.1--4.2 & Post.~4.1--Def.~4.3 \\
Kap.~5: Navigatoren & 5.1--5.5 & Def.~5.1--Def.~5.3 \\
\addlinespace
\multicolumn{3}{@{}l}{\itshape Kapittel 6--9 nedanfor:} \\
\addlinespace
%% -- Kap. 6 --
Post.\ substrat-uavh. (6) & 5.12 & Post.~7.1 \\
Posisjonsinvarians (6) & 1.4, 1.41 & Thm.~7.2 \\
Konj.\ emergent nav. (6) & 5.62 & Conj.~7.3 \\
\addlinespace
%% -- Kap. 7 --
Post.\ stiavh.\ tilp. (7) & 7.1 & Post.~8.1 \\
Stasjonær stig (7) & 7.11 & Thm.~8.2 \\
Stigen endrar landsk. (7) & 7.12 & Thm.~8.3 \\
Spreiingsrelasjonen (7) & 7.2 & Thm.~8.4 \\
\addlinespace
%% -- Kap. 8 --
Geometrisk signatur (8) & 8.1 & Post.~6.1 \\
Informasjonsteoremet (8) & 8.11 & Thm.~6.2 \\
STOLAR-data (8) & 8.12--8.18 & Rem.~6.2 \\
\addlinespace
%% -- Kap. 9 --
Provisorisk kompr. (9) & 7.4 & Thm.~9.1 \\
Varige tradisjonar (9) & 7.5 & Thm.~9.2 \\
Falsifiserbarheit (9) & 7.7 & Thm.~9.3 \\
Refleksiv navigasjon (9) & 7.61 & --- (ikkje i Math.) \\
\end{longtable}
}


%% =====================================================================
%%  APPENDIKS B: NOTASJONSOVERSYN
%% =====================================================================
\chapter{Notasjonsoversyn}
\label{app:notasjon}

{\small
\begin{longtable}{@{}p{28mm}p{62mm}@{}}
\toprule
\textsc{Symbol} & \textsc{Tyding} \\
\midrule
\endfirsthead
\toprule
\textsc{Symbol} & \textsc{Tyding} \\
\midrule
\endhead
\bottomrule
\endfoot

$\MC$ & Formrommet (morphospace) for funksjonell klasse~$C$ \\
$\bm{x}(a)$ & Eigenskapsvektoren til objekt~$a$ \\
$B$ & Busett (okkupert) region i~$\MC$ \\
$O$ & Open region: fysisk mogleg men ikkje realisert \\
$F$ & Forboden region: bryt minst éi avgrensing \\
\addlinespace
$s_i$ & Seleksjonstrykk nr.~$i$ \\
$\bm{s}$ & Seleksjonstrykkvektoren $(s_1, \ldots, s_k)$ \\
$k$ & Antal seleksjonstrykk \\
$\norm{\bm{s}}$ & Samla styrke av seleksjonstrykka \\
\addlinespace
$f(\bm{x})$ & Tilpassingsfunksjonen \\
$\Phi$ & Aggregeringsfunksjonen: $\R^k \to \R$ \\
$\bm{w}$ & Vektvektoren for seleksjonstrykk \\
$\FitLand$ & Tilpassingslandskapet \\
\addlinespace
$\gamma$ & Stig (trajektorie) gjennom formrommet \\
$S[\gamma]$ & Stiavhengig tilpassingsfunksjonale \\
$g$ & Lagrangian for designtrajektoriar \\
$\eta_i$ & Påverknadsfunksjon (landskapsminne) \\
\addlinespace
$\Nav$ & Navigator: trippelet $(G, \mu, \alpha)$ \\
$G$ & Målsett (target region) \\
$\mu$ & Avstandsfunksjon \\
$\alpha$ & Styringsfelt (steering field) \\
$\partial \Nav$ & Kognitiv grenseflate (light cone) \\
$I(\Nav)$ & Intelligens: evna til å unngå lokale optima \\
\addlinespace
$\Ent(\cdot)$ & Shannon-entropi \\
$\MI(\cdot\,;\,\cdot)$ & Gjensidig informasjon \\
$\sigma^2$ & Varians i formfordelinga \\
$\beta$ & Seleksjonsintensitet (invers ``temperatur'') \\
$H$ & Hessian-matrisa ($-\nabla^2 f$) \\
$\lambda_i$ & Eigenverdi nr.~$i$ av $H$ \\
\addlinespace
$\Pareto$ & Pareto-fronten \\
$\Rn$ & $n$-dimensjonalt euklidsk rom \\
$\R$ & Dei reelle tala \\
$\grad$ & Gradient ($\nabla$) \\
$\inner{\cdot}{\cdot}$ & Indre produkt \\
$\norm{\cdot}$ & Norm \\
\end{longtable}
}


%% =====================================================================
%%  APPENDIKS C: FULLSTENDIG ØVINGSNØKKEL
%% =====================================================================
\chapter{Øvingsnøkkel}
\label{app:ovingsnokkel}

Denne nøkkelen gjev hint og retning for kvar øving. Fullstendige
utrekningar er utelate med vilje --- poenget er at studenten gjer
dei sjølv.

\section*{Kapittel 6}

\paragraph{Øving~\ref{ex:6.1}.}
\emph{Hint:} For termostaten er $G$ = settpunkttemperatur,
$\mu$ = $|T_{\text{no}} - T_{\text{sett}}|$, $\alpha$ = slå av/på
varmeelement. For gradient-descent: $G$ = minimum av tapsfunksjon,
$\mu$ = funksjonsverdi, $\alpha$ = negativ gradient.
For bia: $G$ = heksagonal struktur, $\mu$ = avvik frå ideell
geometri, $\alpha$ = justering gjennom voksproduksjon. Alle tre
oppfyller vilkår (a)--(c), men i radikalt ulike substrat.

\paragraph{Øving~\ref{ex:6.2}.}
\emph{Hint:} Posisjonsinvarians seier at \emph{kvar} du er i
formrommet ikkje avheng av \emph{korleis} du kom dit. Men
prosessen avgjer kva \emph{regionar} som er tilgjengelege.
Døme der dei samsvarar: same stolform laga i CNC og for hand.
Døme der dei ikkje kan: ei dobbelt-kurva flate som krev dampbøying
--- ikkje tilgjengeleg for dreiing.

\paragraph{Øving~\ref{ex:6.3}.}
\emph{Hint:} $\phi_{jl}$ er kommunikasjonen mellom
teammedlemmane: møte, skisser, protestar, kompromiss. Det
kollektive $G$ emergerer frå forhandlinga og er ikkje identisk med
nokon einskild persons designintensjon. Bruk den formelle
strukturen frå konjektur~\ref{conj:L.7.3}.

\paragraph{Øving~\ref{ex:6.4}.}
\emph{Hint:} Vel eigenskapar som fangar topologien: antal hol,
massefyldingsgrad, maks-stress/masse-ratio. Same eigenskapar
brukt på begge systemar plasserer dei i same formrom.

\paragraph{Øving~\ref{ex:6.5}.}
\emph{Hint:} Nøkkelen er å skilje mellom ``kan ikkje enno'' og
``kan aldri''. Du treng å kontrollere for reknekapasitet,
sanseoppløysing, og treningstid. Eit godt eksperiment ville gje
ein algoritme \emph{ubegrensa} reknetid og sjå om han framleis
ikkje kan matche den menneskelege navigatoren.

\paragraph{Øving~\ref{ex:6.6}.}
\emph{Hint:} Gode kandidatar: termittbol, maurtegar,
fuglesvermar, open-source-programvare, Wikipedia. For kvart:
identifiser det kollektive målet som ingen einskild deltakar
spesifiserte.


\section*{Kapittel 7}

\paragraph{Øving~\ref{ex:7.1}.}
\emph{Hint:} Posisjonsoptimering: ``kva er den beste stolen?''
Stigoptimering: ``kva er den beste \emph{vegen} til den
stolen?'' I handverket er prosessen viktig fordi han byggjer
ferdigheit, materialforståing, og tradisjonskunnskap.

\paragraph{Øving~\ref{ex:7.2}.}
\emph{Hint:} Same $\bm{x} \in \MC$, ulik $\gamma$. Originalens
stig kryssar frå ein eksisterande attraktor til ein ny region ---
den opnar formrommet. Replika sin stig er kort og monoton innanfor
ein allereie kartlagt attraktor. Ulik $S[\gamma]$, ulik verdi.

\paragraph{Øving~\ref{ex:7.3}.}
\emph{Hint:} $\frac{\partial g}{\partial \gamma}$ trekkjer mot
høg tilpassing (``gå dit det er bra''). $\frac{d}{dt}
\frac{\partial g}{\partial \dot{\gamma}}$ straffar brå retningsendring
(``ikkje hopp''). Fyrste ledd dominerer i tidleg
utforsking; andre ledd dominerer i moden perfeksjonering.

\paragraph{Øving~\ref{ex:7.4}.}
\emph{Hint:} Vel t.d.\ medisinsk utstyr (sterkt) og
dekorative vegglamper (svakt). Prediker at variansen i
geometri for vegglamper er mykje høgare enn for medisinsk utstyr.
Test ved å måle kompaktheit og kurvatur for eit representativt
utval og samanlikne spreiinga.

\paragraph{Øving~\ref{ex:7.5}.}
\emph{Hint:} Bruk eit 2D-formrom, to tradisjonar som startar i
same punkt men tek ulike ruter. Definer
$\eta_i(\bm{y}, \bm{x})$ som ein enkel funksjon (t.d.\ ein
Gauss-klokke rundt $\bm{y}$). Vis at tradisjon~A sin rute
forsterkar ein annan topp enn tradisjon~B sin rute.

\paragraph{Øving~\ref{ex:7.6}.}
\emph{Hint:} I eit éindimensjonalt rom er formsekvensen ei linje.
I $n$~dimensjonar kan stigen forgreine seg, sløyfe tilbake, og
krysse seg sjølv. Kublers typologi (primærobjekt, replikar)
vert meir kompleks: eit objekt kan vere primært langs éin
dimensjon og ein replika langs ein annan.


\section*{Kapittel 8}

\paragraph{Øving~\ref{ex:8.1}.}
\emph{Hint:} Innanfor ein funksjonell klasse er funksjonen same
for alle objekt. Å vite funksjonen (``det er ein stol'') gjev deg
null ny informasjon om forma, fordi du allereie visste det.

\paragraph{Øving~\ref{ex:8.2}.}
\emph{Hint:} Plast er lett å forme i komplekse kurver
(sprøytestøyping), men har avgrensa strukturell stivleik. Forventa
profil: høg kompaktheit (lukka former), middels symmetri, høg
kurvatur. Grunngjev med dei fysiske eigenskapane til sprøytestøypt
polypropylen.

\paragraph{Øving~\ref{ex:8.3}.}
\emph{Hint:} Med plast lik fordelt (5/5) er plast uavhengig av
geometri. Den nye $\MI$ vert lågare enn 0{,}19~bit fordi den
tredje kategorien ``vatnar ut'' sambandet. Rekn ut med same
formel; den nye tabellen er $3 \times 2$.

\paragraph{Øving~\ref{ex:8.4}.}
\emph{Hint:} $\norm{\bm{s}}$ er den samla styrken av
seleksjonstrykka. I perioden 1750--1850 vart eitt trykk
(materialpresteige) svært sterkt medan andre trykk (lokal
tradisjon, individuell innovasjon) vart undertrykte av kolonialt
handelsnettverk. Resultat: $\norm{\bm{s}}$ veks,
$\sigma^2$ krympar, entropien kollapsar.

\paragraph{Øving~\ref{ex:8.5}.}
\emph{Hint:} For knivar: geometriske eigenskapar = bladlengd,
bladbreidde, kurvatur, vekt. Materialkategoriar = karbonstål,
rustfritt stål, keramikk. Kontroller for brukstype (kokkekniv vs
jaktkniv) for å halde funksjonell klasse konstant.
Bruk kji-kvadrat-test eller permutasjonstest på den betinga
fordelinga.

\paragraph{Øving~\ref{ex:8.6}.}
\emph{Hint:} $P(\bm{x} \mid m) = P(\bm{x})$ krev at materialet
har \emph{null} innverknad på forma. Det ville krevje at alle
material har identiske fysiske eigenskapar (same brotstyrke, same
bøyelighet, same bearbeidingskarakter) --- i praksis umogleg.
Alternativt ville det krevje ein formgjevar som fullstendig
ignorerer materialets respons, noko som bryt med den taktile
tilbakemeldingssløyfa i all fysisk formgjeving.


\section*{Kapittel 9}

\paragraph{Øving~\ref{ex:9.1}.}
\emph{Hint:} Provisorisk kompromiss = optimalt akkurat \emph{no},
men ikkje for alltid. Døme: ein fastlinetelefon var optimal for
telekommunikasjon i 1960, men er suboptimal i 2026 fordi
seleksjonstrykka (mobilitet, bandbreidd, forventningar) har
endra seg drastisk.

\paragraph{Øving~\ref{ex:9.2}.}
\emph{Hint:} Eit ufalsifiserbart utsegn kan ikkje testast og gjev
dermed ikkje vitskapleg kunnskap. ``Form følgjer funksjon'' var
ufalsifiserbart --- det forklarte alt og dermed ingenting.
Falsifiserbarheit tyder at teorien \emph{risikerer} å ta feil,
og det er nettopp denne risikoen som gjev han forklaringskraft.

\paragraph{Øving~\ref{ex:9.3}.}
\emph{Hint:} Følg strukturen frå seksjon~9.5. For ei lampe laga
for hand: (a)~handverkar + verktøy, (b)~former i tre/metall/glas
innanfor verktøyets avgrensingar, (c)~lyskvalitet, pris, estetikk,
(d)~ein multimodal overflate med toppar rundt klassiske lampetype-ar,
(e)~eksplorasjon = prototypar, eksploitering = perfeksjonering av
ein utvald form.

\paragraph{Øving~\ref{ex:9.4}.}
\emph{Hint:} Gode kandidatar: japansk sverdsmiing, norsk
rosemåling, italiensk glasblåsing. For kvar: identifiser
materialforståing, stilistisk tradisjon, individuell meisterskap,
institusjonell overføring (meister--lærling), og marknads\-tilpassing
som separate mekanismar. Vis at kvar mekanisme responderer på eit
anna trykk.

\paragraph{Øving~\ref{ex:9.5}.}
\emph{Hint:} Dersom $\Phi$ er ein funksjon av
$\gamma([0, t])$, er tilpassingsfunksjonen sjølv stiavhengig.
Euler--Lagrange-likninga vert meir kompleks: lagrangianen $g$
avheng ikkje berre av $\gamma(t)$ og $\dot{\gamma}(t)$, men av
heile den tidlegare stigen. Dette gjev ein \emph{integrodifferensiallikning}
i staden for ein ordinær differensiallikning.

\paragraph{Øving~\ref{ex:9.6}.}
\emph{Hint:} Generer $N$~bilete av ``stol'' ved temperatur
$T_1 = 0{,}3$, $T_2 = 0{,}7$ og $T_3 = 1{,}0$. Mål
kompaktheit og symmetri for kvart bilete. Rekn ut variansen for
kvar temperaturnivå. Spreiingsrelasjonen predikerer
$\sigma^2 \propto T$ (sidan $T = 1/\beta$). Plott
$\sigma^2$ mot $T$ og sjekk om samanhengen er tilnærma lineær.

\paragraph{Øving~\ref{ex:9.7}.}
\emph{Hint:} For ``robustheit for generativ AI = låg'': argumentet
\emph{for} er at nettverket er trengt på ei fast datamengd og kan
ikkje tilpasse seg nye seleksjonstrykk utan re-trening.
Argumentet \emph{mot} er at fine-tuning og RLHF gjev ein viss
adaptiv kapasitet. Data du treng: mål kor raskt ein
generativ modell produserer suboptimale resultat etter eit dokumentert
skifte i brukarpreferansar, samanlikna med kor raskt ein
handverkar tilpassar seg.


%% ══════════════════════════════════════════════════════════════════
%% APPENDIKS
%% ══════════════════════════════════════════════════════════════════

\appendix

\chapter{Korrespondansetabell}

Tabellen under viser korleis kvart kapittel i denne læreboka svarar
til proposisjonar i \textit{Tractatus Formae} og formelle element i
\textit{Forml\ae re: Mathematica}.

\medskip
\begin{small}
\begin{longtable}{@{}p{38mm}p{20mm}p{25mm}@{}}
\toprule
\textsc{L\ae rebok} & \textsc{Tractatus} & \textsc{Mathematica} \\
\midrule
\endfirsthead
\toprule
\textsc{L\ae rebok} & \textsc{Tractatus} & \textsc{Mathematica} \\
\midrule
\endhead
\bottomrule
\endfoot
Kap.~1 Formrommet   & 1.1--1.6  & Def.~1.1--1.2, Prop.~1.4 \\
Kap.~2 Seleksjonstrykk & 2.01--2.6 & Def.~2.1--2.3, Post.~2.4, 2.6, Thm.~2.5, 2.7, 2.8 \\
Kap.~3 Tilpassingslandskapet & 3.1--3.82 & Def.~3.1--3.2, 3.4--3.6, Thm.~3.3, Cor.~3.5 \\
Kap.~4 Dynamikk     & 4.1--4.22 & Post.~4.1, Thm.~4.2, Def.~4.3 \\
Kap.~5 Navigatorar  & 5.1--5.6  & Def.~5.1--5.3 \\
Kap.~6 Substrat og emergens & 5.12--5.62 & Post.~7.1, Thm.~7.2, Conj.~7.3 \\
Kap.~7 Variasjonsprinsipp & 7.1--7.3 & Post.~8.1, Thm.~8.2--8.4 \\
Kap.~8 Materialsignaturar & 8.1--8.18 & Post.~6.1, Thm.~6.2 \\
Kap.~9 Konsekvensar & 7.4--7.7  & Thm.~9.1--9.3 \\
\end{longtable}
\end{small}


\chapter{Notasjonsoversyn}

\begin{description}[style=nextline, leftmargin=2em]
  \item[$\MC$] Formrommet til klasse $C$ (kompakt delmengd av $\Rn$)
  \item[$s_i$] Seleksjonstrykk nr.~$i$ ($s_i : \MC \to \R$)
  \item[$\bm{s}$] Seleksjonstrykkvektoren $(s_1, \ldots, s_k)$
  \item[$f$] Tilpassingsfunksjonen $f = \Phi(\bm{s})$
  \item[$\FitLand$] Tilpassingslandskapet (grafen til $f$)
  \item[$\Nav$] Navigator-trippelet $(G, \mu, \alpha)$
  \item[$\CogBound$] Kognitiv grenseflate (lyskjegle)
  \item[$B, O, F$] Busett, ope og forbode region
  \item[$\gamma$] Stig (vandring) gjennom formrommet
  \item[{$S[\gamma]$}] Stiavhengig tilpassingsverdi (verknadsfunksjonalet)
  \item[$\Ent$] Shannon-entropi
  \item[$\MI$] Gjensidig informasjon
  \item[$\Pareto$] Pareto-fronten
\end{description}
